\chapter{Herpangina}
\index{Herpangina}
\label{sec:Herpangina}

\section{Overview}
Herpangina is a common childhood viral illness characterized by fever and painful vesicular lesions on the soft palate, tonsils, and posterior pharynx. It is typically caused by Coxsackie A viruses (group A enteroviruses) and most commonly affects children aged 3--10 years. The condition is self-limited, with symptoms resolving within 7--10 days.
\section{Classification}

\begin{description}[font=\normalfont\bfseries,leftmargin=3em,labelwidth=2.5em]
  \item[Cause] Infectious (Viral)
  \item[Organ System] Pediatric / Infectious
  \item[Pathophysiology] Enterovirus infection causing vesicular lesions in the posterior oropharynx
  \item[Duration] Acute (7--10 days)
  \item[Onset] Sudden
  \item[Mode of Transmission] Fecal-oral and respiratory droplets
  \item[Clinical Course] Acute, self-limited
  \item[Severity] Mild to Moderate
  \item[Extent of Disease] Localized (oropharynx) with systemic symptoms
  \item[Age Group] Children (3--10 years)
  \item[Epidemiology] Common; epidemics in summer and early fall in temperate climates
  \item[ICD-11 Code] 1D64.2
\end{description}

\section{Causes}
Herpangina is caused by group A enteroviruses, most commonly Coxsackie A viruses (serotypes 1--10, 16, 22), and less commonly by Coxsackie B viruses, echoviruses, and enterovirus 71. The virus enters through the gastrointestinal tract, replicates in lymphoid tissues, and causes viremia with subsequent infection of the oropharyngeal mucosa, leading to vesicle formation.

\section{Risk Factors}

\begin{itemize}[noitemsep]
  \item Age 3--10 years (peak incidence)
  \item Summer and early fall season
  \item Daycare or school attendance
  \item Poor hand hygiene
  \item Household contact with infected individuals
  \item Lack of prior immunity
\end{itemize}

\section{Signs and Symptoms}

\subsection{Early symptoms}
\begin{itemize}[noitemsep]
  \item Early manifestations of herpangina may be subtle and nonspecific
\end{itemize}

\subsection{Common symptoms}
\begin{itemize}[noitemsep]
  \item \subsection{Early symptoms}
  \item \subsection{Common symptoms}
  \item \textbf{Common symptoms:}
  \item Sudden onset of high fever (38--40°C)
  \item Sore throat and pain with swallowing (odynophagia)
  \item Small (1--2 mm) vesicles on the soft palate, uvula, tonsils, and posterior pharynx
  \item Vesicles progress to shallow ulcers with a gray base and erythematous halo
  \item Typically 2--6 lesions (fewer than hand, foot, and mouth disease)
  \item Headache, malaise, anorexia, and vomiting
  \item Pain or discomfort in the affected area
  \item Difficulty performing daily activities
\end{itemize}

\subsection{Less common symptoms}
\begin{itemize}[noitemsep]
  \item Atypical or less common presentations of herpangina may occur
\end{itemize}

\subsection{Emergency warning signs}
\begin{itemize}[noitemsep]
  \item Severe or worsening symptoms of herpangina require immediate medical evaluation
\end{itemize}
\section{Progression and Stages}
After an incubation period of 3--6 days, fever and sore throat develop abruptly. Vesicles appear within 24--48 hours and progress to shallow ulcers within 1--2 days. The ulcers heal spontaneously within 7--10 days without scarring. Fever typically resolves within 2--4 days.

\section{Complications}

\begin{itemize}[noitemsep]
  \item Dehydration from painful swallowing and decreased oral intake
  \item Febrile seizures (uncommon)
  \item Viral meningitis (rare, particularly with enterovirus 71)
  \item Weight loss from decreased oral intake
  \item Reduced quality of life
  \item Emotional and psychological effects
\end{itemize}

\section{Diagnosis}

\begin{itemize}[noitemsep]
  \item Clinical diagnosis based on characteristic lesion appearance and distribution
  \item Throat swab for viral culture or PCR for confirmation
  \item Enterovirus serology (rarely needed)
  \item Complete blood count may show lymphocytosis
  \item Distinguish from hand, foot, and mouth disease (lesions on hands and feet)
\end{itemize}

\section{Prevention}

\begin{itemize}[noitemsep]
  \item Good hand hygiene, especially after diaper changes
  \item Disinfection of contaminated surfaces and toys
  \item Avoid sharing cups, utensils, and towels
  \item Exclusion from school or daycare while febrile
  \item Go for regular check-ups so any problems are caught early
\end{itemize}

\section{Treatment Options}

\begin{itemize}[noitemsep]
  \item Supportive care: antipyretics for fever (acetaminophen, ibuprofen)
  \item Adequate hydration (cold fluids, popsicles, ice cream for comfort)
  \item Topical analgesics for oral lesions (lidocaine mouthwash, viscous lidocaine)
  \item Avoid acidic or salty foods that irritate ulcers
  \item Monitor for dehydration and provide IV fluids if needed
  \item Lifestyle changes and self-care
\end{itemize}

\section{Living with It}

\begin{itemize}[noitemsep]
  \item Follow your treatment plan as prescribed
  \item Keep up with regular appointments with your healthcare provider
  \item Adjust your daily habits as needed
  \item Reach out to family, friends, or support groups for help
  \item Keep learning about your condition and how to manage it
\end{itemize}

\section{Outlook}
Herpangina is a self-limited illness with an excellent prognosis. Symptoms typically resolve completely within 7--10 days without specific antiviral therapy. The main concern is dehydration from painful swallowing, which requires careful attention to fluid intake. Complications are rare in immunocompetent children. Immunity is serotype-specific, so children can have multiple episodes.
